% Shared preamble for the Generative AI (CS 782/882, Fall 2026) lecture notes.
% Each weekNN-notes.tex file inlines this block so it compiles on its own.
\documentclass[11pt]{article}
\usepackage[margin=1in]{geometry}
\usepackage[T1]{fontenc}
\usepackage[utf8]{inputenc}
\usepackage{lmodern}
\usepackage{amsmath,amssymb,amsthm,mathtools,bm}
\usepackage{enumitem}
\usepackage{graphicx}
\usepackage{booktabs}
\usepackage{xcolor}
\usepackage{tikz}
\usetikzlibrary{arrows.meta,positioning,calc,shapes}
\usepackage{listings}
\usepackage[colorlinks=true,linkcolor=blue!60!black,urlcolor=blue!60!black,citecolor=blue!60!black]{hyperref}

\newcommand{\E}{\mathbb{E}}
\newcommand{\Var}{\mathrm{Var}}
\newcommand{\KL}{\mathrm{KL}}
\newcommand{\given}{\,\vert\,}
\newcommand{\R}{\mathbb{R}}
\DeclareMathOperator*{\argmax}{arg\,max}
\DeclareMathOperator*{\argmin}{arg\,min}
\newtheorem{theorem}{Theorem}
\newtheorem{lemma}{Lemma}
\newtheorem{definition}{Definition}
\theoremstyle{remark}
\newtheorem*{remark}{Remark}

\lstset{basicstyle=\ttfamily\small,breaklines=true,frame=single,columns=fullflexible}

\title{Supplementary Notes: Conformal Prediction}
\author{Hong Qin, Old Dominion University}
\date{Fall 2026}

\begin{document}
\maketitle
\begin{center}\small\itshape Note: this topic was covered in the sixteen-week version of this course. It is provided here as optional reading and is not part of the eight-week core. Nothing in these notes is required for the weekly assignments or the course project.\end{center}
\medskip

\section{Introduction}
Conformal prediction (CP) is a framework for \textbf{uncertainty quantification (UQ)} in machine learning.
It wraps around any black-box model to produce prediction \emph{sets} or \emph{intervals} with rigorous, distribution-free coverage guarantees.
Unlike Bayesian or ensemble approaches, CP does not assume any specific data distribution, only that the data are \textbf{exchangeable} (i.i.d.\ is a common special case).

\bigskip
The central guarantee is that for a user-chosen error rate $\alpha \in (0,1)$,
\[
\mathbb{P}\!\left( Y_{\text{true}} \notin \hat{C}(X) \right) \leq \alpha,
\]
where $\hat{C}(X)$ is the prediction set for input $X$. In words: with probability at least $1-\alpha$, the prediction set will contain the true label or value.

\section{Core Principle and Intuition}
The idea is to measure how ``typical'' or ``atypical'' a new example looks compared to past data.
Instead of producing a single prediction (which may be over- or under-confident), CP produces a set of possible outcomes that is guaranteed to contain the truth at least $1-\alpha$ of the time.

\bigskip
\noindent \textbf{Key steps:}
\begin{enumerate}
    \item Split data into a training set (for fitting the base model) and a calibration set (for calibration).
    \item Train any base model $f$ (regressor, classifier, neural network, etc.).
    \item Define a \textbf{nonconformity score} $s(x,y)$, which measures how unusual $(x,y)$ is relative to the model.
    \item Compute nonconformity scores for all calibration points.
    \item Take a high quantile of these scores as the threshold.
    \item For a new input, include in the prediction set all labels/values whose score is below the threshold.
\end{enumerate}

\section{Nonconformity vs.\ Similarity Scores}
A central concept is the \textbf{nonconformity score}, which quantifies how ``strange'' or ``non-conforming'' a point looks.

\begin{itemize}
    \item Regression: $s(x,y) = |y - f(x)|$ (absolute residual).
    \item Classification: $s(x,y) = 1 - \hat{p}(y|x)$ (one minus predicted probability of the true class).
\end{itemize}

Why not use similarity?
If we used a similarity score (high = typical, low = strange), we would need to invert all comparisons to identify outliers.
Using nonconformity is natural: larger scores automatically correspond to stranger examples.
This makes it straightforward to control the ``tail risk'' by picking a quantile threshold.

\bigskip
\noindent \textbf{Analogy:}
Grading exam answers:
\begin{itemize}
    \item Similarity view: ``This answer is 90\% similar to the key.''
    \item Nonconformity view: ``This answer makes 10\% mistakes.''
\end{itemize}
To set a rule like ``reject the worst 5\% of answers,'' it is more natural to work with mistakes (nonconformity).

\section{Role of $\alpha$}
The parameter $\alpha$ is the \textbf{user-chosen error rate}:
\[
\mathbb{P}\!\left( Y_{\text{true}} \in \hat{C}(X) \right) \geq 1-\alpha.
\]

\begin{itemize}
    \item Smaller $\alpha$ (e.g.\ 0.05) means you want to be wrong only 5\% of the time. Prediction sets are larger/wider.
    \item Larger $\alpha$ (e.g.\ 0.2) means you accept being wrong 20\% of the time. Prediction sets are smaller/tighter.
\end{itemize}

Thus, $\alpha$ directly controls the \textbf{tradeoff between coverage and efficiency}.
It is not the same as a Bayesian posterior probability; instead it is a distribution-free guarantee.

\section{Quantile Selection in Practice}
To achieve the target error rate $\alpha$, we use the nonconformity scores from the calibration set:
\begin{enumerate}
    \item Suppose calibration set size is $n$.
    \item Sort the scores: $s_{(1)} \leq s_{(2)} \leq \dots \leq s_{(n)}$.
    \item Choose threshold $q = s_{(k)}$ where
    \[
    k = \left\lceil (1-\alpha)(n+1) \right\rceil.
    \]
\end{enumerate}

This ensures finite-sample validity.

\subsection*{Regression Example}
Calibration residuals:
\[
\{0.5, 1.2, 0.9, 2.1, 1.8, 0.3, 1.1, 2.5, 0.7, 1.4\}.
\]
Sorted:
\[
\{0.3, 0.5, 0.7, 0.9, 1.1, 1.2, 1.4, 1.8, 2.1, 2.5\}.
\]

With $\alpha = 0.1$ and $n=10$:
\[
k = \lceil 0.9 \times 11 \rceil = \lceil 9.9 \rceil = 10.
\]
Threshold $q=2.5$.
If $\hat{y}=8.0$, prediction interval is $[8.0-2.5,\,8.0+2.5]=[5.5,\,10.5]$.

\subsection*{Classification Example}
Suppose calibration predicted probabilities for the true class are
\[
\{0.90, 0.70, 0.60, 0.50, 0.40, 0.30\}.
\]
Scores (1 minus):
\[
\{0.10, 0.30, 0.40, 0.50, 0.60, 0.70\}.
\]
Sorted: same order.

With $\alpha=0.1$, $n=6$:
\[
k = \lceil 0.9 \times 7 \rceil = \lceil 6.3 \rceil = 7,
\]
so we take the largest score $0.70$.

For a new input with predicted probabilities
\[
p(A)=0.65,\; p(B)=0.25,\; p(C)=0.10,
\]
scores are $0.35, 0.75, 0.90$.
Keep labels with score $\leq 0.70$. Prediction set = $\{A\}$.

\bigskip
If instead $p(A)=0.40,\; p(B)=0.35,\; p(C)=0.25$,
scores are $0.60, 0.65, 0.75$.
Prediction set = $\{A,B\}$.

\section{Key Takeaways}
\begin{itemize}
    \item Conformal prediction produces set-valued predictions with distribution-free coverage guarantees.
    \item Nonconformity scores measure ``strangeness'' and make it easy to threshold extreme cases.
    \item The parameter $\alpha$ is the miscoverage tolerance; $1-\alpha$ is the coverage guarantee.
    \item Prediction sets adapt to model quality: narrow when the model is accurate, wide when the model is uncertain.
\end{itemize}

\section{Conformal Prediction for Time Series}

While conformal prediction provides distribution-free coverage guarantees under exchangeability,
time series data violate this assumption due to temporal dependence and potential non-stationarity.
Extending CP to time series requires careful handling of these dependencies.

\subsection*{Challenges}
\begin{itemize}
    \item \textbf{Dependence:} Observations at different time points are correlated, so the usual permutation-based arguments for validity no longer hold directly.
    \item \textbf{Non-stationarity:} Time series may drift or exhibit changing variance, requiring adaptive calibration.
    \item \textbf{Prediction horizon:} For $h$-step ahead forecasts, uncertainty grows with horizon length.
\end{itemize}

\subsection*{Approaches}
Several strategies have been proposed to adapt CP to time series:
\begin{itemize}
    \item \textbf{Sliding-window conformal prediction:} Use only the most recent $w$ observations as the calibration set. This enforces a form of ``local exchangeability'' and adapts to drift.
    \item \textbf{Block conformal prediction:} Partition the series into blocks and treat each block as approximately exchangeable, capturing short-range dependencies.
    \item \textbf{Adaptive/weighted conformal prediction:} Weight calibration residuals by recency or importance, giving more influence to recent dynamics.
    \item \textbf{Model-based CP:} Combine CP with probabilistic forecasters (e.g., autoregressive models, RNNs, Transformers) to generate residual-based nonconformity scores that account for temporal structure.
\end{itemize}

\subsection*{Example: Sliding-window CP}
Suppose we want to forecast $Y_{t+1}$ from features $X_t$.
We train a predictive model $f$ on past data, then compute nonconformity scores on a moving calibration window:
\[
s_i = |Y_i - f(X_{i-1})|, \quad i = t-w+1, \dots, t.
\]
We then compute the $(1-\alpha)$-quantile of $\{s_i\}$ and form the prediction interval
\[
\hat{C}_{t+1} = \big[ f(X_t) - q_{1-\alpha}, \, f(X_t) + q_{1-\alpha} \big].
\]
As $t$ advances, the calibration window updates, adapting to non-stationary trends.

\subsection*{Applications}
Recent work has successfully applied conformal prediction to time series forecasting in domains such as:
\begin{itemize}
    \item \textbf{Traffic forecasting:} Providing uncertainty bands around flow predictions on road networks.
    \item \textbf{Energy demand forecasting:} Quantifying risk in power grid load prediction.
    \item \textbf{Finance:} Forecasting with distribution-free confidence intervals under non-stationarity.
\end{itemize}

\subsection*{Related Work}
\begin{itemize}
    \item C. Xu and Y. Xie. ``Conformal Prediction Interval for Dynamic Time-Series.'' In \emph{Proceedings of the 38th International Conference on Machine Learning (ICML)}, 2021.
    \item I. Gibbs and E. Cand\`es. ``Adaptive Conformal Inference Under Distribution Shift.'' In \emph{Advances in Neural Information Processing Systems (NeurIPS)}, 2021.
    \item A. N. Angelopoulos and S. Bates. ``A Gentle Introduction to Conformal Prediction and Distribution-Free Uncertainty Quantification.'' arXiv:2107.07511, 2021.
\end{itemize}

These works extend CP to dependent data by developing rigorous guarantees under relaxed assumptions (e.g., mixing conditions, block exchangeability) and by designing adaptive calibration strategies.

\end{document}
